\chapter[Bethe Ansatz and Yang--Baxter Integrability]{Preview of Bethe Ansatz and Yang--Baxter Integrability}
\label{chap:bethe-yang-baxter-preview}

The preceding chapters developed two main solution mechanisms.  First, transfer
matrices convert a classical statistical model into a spectral problem.  Second,
Jordan--Wigner and Bogoliubov transformations solve a large class of spin chains
because their fermionic images are quadratic.  The next integrability mechanism
is different.  Bethe-ansatz integrability solves genuinely interacting models,
not by reducing them to independent one-particle modes, but by exploiting the
factorization of many-body scattering and the existence of an infinite commuting
family of transfer matrices.

This chapter is a conceptual and notational preview.  Its purpose is to explain
what the Bethe ansatz solves, why the Yang--Baxter equation appears, and how the
row-to-row transfer matrices of vertex models become quantum Hamiltonians such as
the XXZ chain.  The complete coordinate and algebraic Bethe-ansatz derivations
are postponed to later chapters.

\section{From free particles to factorized scattering}
\label{sec:ba-factorized-scattering}

A free-fermion Hamiltonian is solvable because it can be diagonalized into
independent quasiparticles,
\begin{equation}
    H=E_0+\sum_k E_k\gamma_k^\dagger\gamma_k .
\end{equation}
The many-body spectrum is then obtained by filling one-particle levels.  Bethe
ansatz uses a weaker but still powerful property: particles interact, but the
scattering is elastic and factorized.

\begin{definition}[Factorized scattering]
Consider a one-dimensional many-body system with stable asymptotic quasiparticles
labelled by momenta or rapidities.  The scattering is called \emph{factorized} if
an $M$-body collision can be decomposed into an ordered product of two-body
scattering processes, with no particle production and no diffraction.
\end{definition}

In one spatial dimension the ordering of particles is meaningful.  For
coordinates
\begin{equation}
    x_1<x_2<\cdots<x_M,
\end{equation}
a coordinate Bethe wave function takes the form
\begin{equation}
\label{eq:coordinate-ba-wavefunction-preview}
    \Psi(x_1,\ldots,x_M)
    =
    \sum_{P\in S_M} A(P)
    \exp\left(\ii\sum_{a=1}^{M}k_{P_a}x_a\right),
\end{equation}
where $P$ permutes the momenta among the ordered particles.  Interactions do not
change the set of momenta.  They only relate amplitudes associated with different
permutations.  For neighboring exchanges one writes schematically
\begin{equation}
\label{eq:two-body-scattering-amplitude-preview}
    A(P\circ s_a)
    =S(k_{P_a},k_{P_{a+1}})A(P),
\end{equation}
where $s_a$ exchanges the $a$th and $(a+1)$st momenta.  For a scalar scattering
phase, factorization is almost automatic because complex numbers commute.  For
particles with internal spin, color, or auxiliary indices, $S$ becomes a matrix,
and consistency is highly nontrivial.

The basic physical meaning is the following.  In a generic interacting model, a
many-particle collision can redistribute momenta in ways that cannot be described
as a sequence of independent two-body events.  In an integrable one-dimensional
model, the infinite family of conservation laws is strong enough to forbid such
diffraction.  The only allowed effect of scattering is to permute the rapidities
and multiply the wave function by two-body scattering matrices.

\begin{principle}[Bethe-ansatz solvability]
A one-dimensional interacting model is Bethe-ansatz solvable when its eigenstates
can be assembled from plane waves whose amplitudes are related by consistent
factorized two-body scattering.
\end{principle}

\section{A two-magnon preview}
\label{sec:two-magnon-preview}

The simplest way to see the Bethe idea is to consider a spin chain above a
ferromagnetic reference state.  Let
\begin{equation}
    \ket{0}=\ket{\uparrow\uparrow\cdots\uparrow}
\end{equation}
and let a down spin be interpreted as a magnon.  A one-magnon state is
\begin{equation}
    \ket{k}=\sum_{x=1}^{L}\ee^{\ii kx}\sigma_x^-\ket{0} .
\end{equation}
For a translation-invariant chain, it is an eigenstate in the one-magnon sector.
The nontrivial part begins in the two-magnon sector, where the wave function has
one expression for $x_1<x_2$:
\begin{equation}
\label{eq:two-magnon-wavefunction-preview}
    \Psi(x_1,x_2)
    =A_{12}\ee^{\ii(k_1x_1+k_2x_2)}
    +A_{21}\ee^{\ii(k_2x_1+k_1x_2)} .
\end{equation}
The interaction only matters when the two magnons approach each other.  The
Schrodinger equation at contact fixes the ratio
\begin{equation}
\label{eq:two-magnon-S-preview}
    A_{21}=S(k_1,k_2)A_{12} .
\end{equation}
For a periodic chain, moving one magnon around the ring makes it scatter once
with every other magnon.  This gives the schematic Bethe equations
\begin{equation}
\label{eq:schematic-bethe-equations-preview}
    \ee^{\ii k_j L}
    =
    \prod_{\ell\neq j}S(k_j,k_\ell),
    \qquad j=1,\ldots,M .
\end{equation}
Thus periodic boundary conditions quantize the momenta in an interaction-dependent
way.  Free particles would have $S=\pm1$ and hence ordinary integer or
half-integer momenta.  Interacting Bethe-ansatz particles acquire phase shifts
from all other particles.

\begin{example}[The XXZ two-body scattering phase]
For one common Pauli-matrix convention, the XXZ chain may be written as
\begin{equation}
\label{eq:xxz-preview-pauli-convention}
    H_{\rm XXZ}
    =
    J\sum_{j=1}^{L}
    \left(
    \sigma_j^x\sigma_{j+1}^x
    +\sigma_j^y\sigma_{j+1}^y
    +\Delta(\sigma_j^z\sigma_{j+1}^z-1)
    \right).
\end{equation}
In the two-magnon coordinate Bethe ansatz, the contact equation gives a scattering
phase of the form
\begin{equation}
\label{eq:xxz-momentum-S-preview}
    S(k_1,k_2)
    =-
    \frac{1+\ee^{\ii(k_1+k_2)}-2\Delta\ee^{\ii k_1}}
         {1+\ee^{\ii(k_1+k_2)}-2\Delta\ee^{\ii k_2}} .
\end{equation}
The precise overall sign and energy convention depend on the normalization of
\cref{eq:xxz-preview-pauli-convention}, but the structural message is universal:
the interaction is encoded in a two-body phase shift, and the many-body momenta
are fixed by \cref{eq:schematic-bethe-equations-preview}.
\end{example}

It is often more convenient to replace momenta by rapidities.  In the XXX limit,
the Bethe equations are commonly written as
\begin{equation}
\label{eq:xxx-bethe-equations-preview}
    \left(
    \frac{\lambda_j+\ii/2}{\lambda_j-\ii/2}
    \right)^L
    =
    \prod_{\ell\neq j}
    \frac{\lambda_j-\lambda_\ell+\ii}
         {\lambda_j-\lambda_\ell-\ii} .
\end{equation}
For the gapless XXZ chain with $\Delta=\cos\gamma$, a standard rapidity form is
\begin{equation}
\label{eq:xxz-bethe-equations-preview}
    \left[
    \frac{\sinh(\lambda_j+\ii\gamma/2)}
         {\sinh(\lambda_j-\ii\gamma/2)}
    \right]^L
    =
    \prod_{\ell\neq j}
    \frac{\sinh(\lambda_j-\lambda_\ell+\ii\gamma)}
         {\sinh(\lambda_j-\lambda_\ell-\ii\gamma)} .
\end{equation}
The left-hand side is free propagation around the ring.  The right-hand side is
the accumulated product of two-body scattering phases.

\section{Why a three-body consistency condition is needed}
\label{sec:three-body-consistency-preview}

For three particles there are multiple ways to exchange the same initial and
final orderings.  For instance, the permutation reversing $(1,2,3)$ to $(3,2,1)$
can be decomposed as
\begin{equation}
    s_1s_2s_1=s_2s_1s_2 .
\end{equation}
If the scattering amplitudes are scalar phases, this identity causes no serious
constraint.  If the two-body scatterings are matrices acting on internal degrees
of freedom, the two possible sequences must give the same operator.  This is the
origin of the Yang--Baxter equation.

Let $R_{ab}(u)$ denote a two-body scattering or vertex operator acting nontrivially
on spaces $a$ and $b$, with spectral parameter difference $u$.  The Yang--Baxter
equation is
\begin{equation}
\label{eq:yang-baxter-equation-preview}
    R_{12}(u-v)R_{13}(u-w)R_{23}(v-w)
    =
    R_{23}(v-w)R_{13}(u-w)R_{12}(u-v) .
\end{equation}
It is an equality of operators on a triple tensor product.  Diagrammatically it
states that two different orders of resolving a three-line crossing are
equivalent.  Algebraically it is precisely the consistency condition required for
factorized scattering with internal degrees of freedom.

\begin{remark}[Yang--Baxter versus ordinary commutativity]
The Yang--Baxter equation is not the statement that two matrices commute.  It is
a cubic relation among three two-body operators acting on three tensor factors.
Its consequence, after constructing a monodromy matrix and tracing over an
auxiliary space, is the commutativity of transfer matrices.
\end{remark}

\section{R-matrix as vertex weight and as scattering matrix}
\label{sec:r-matrix-two-interpretations}

The same $R$-matrix has two complementary meanings.

First, in a two-dimensional classical vertex model, $R$ assigns a local Boltzmann
weight to a vertex.  The two tensor indices may be viewed as spin or arrow states
living on the lines entering and leaving the vertex.  The partition function is
built by multiplying local $R$-matrix weights and summing over internal indices.

Second, in a one-dimensional quantum chain, $R$ is the algebraic object that
controls two-body scattering and the commuting transfer matrices.  The spectral
parameter $u$ plays the role of an auxiliary rapidity.  Taking a logarithmic
derivative of the transfer matrix produces a quantum Hamiltonian.

For the six-vertex model, a standard trigonometric or hyperbolic $R$-matrix is
\begin{equation}
\label{eq:six-vertex-R-preview}
    R(u)=
    \begin{pmatrix}
    a(u)&0&0&0\\
    0&b(u)&c(u)&0\\
    0&c(u)&b(u)&0\\
    0&0&0&a(u)
    \end{pmatrix},
\end{equation}
where, for the hyperbolic parametrization,
\begin{equation}
\label{eq:six-vertex-weights-hyperbolic-preview}
    a(u)=\sinh(u+\eta),
    \qquad
    b(u)=\sinh u,
    \qquad
    c(u)=\sinh\eta .
\end{equation}
The nonzero entries encode the ice rule: the number of incoming arrows equals the
number of outgoing arrows at each vertex.  In the associated XXZ chain, the
anisotropy parameter is controlled by
\begin{equation}
\label{eq:xxz-anisotropy-preview}
    \Delta=\cosh\eta
\end{equation}
for this hyperbolic convention.  In the gapless convention one uses sines instead
of hyperbolic sines and obtains $\Delta=\cos\gamma$.

\section{Monodromy matrix and commuting transfer matrices}
\label{sec:monodromy-transfer-preview}

To build a length-$L$ quantum chain, introduce an auxiliary space $a$ and quantum
spaces $1,\ldots,L$.  The monodromy matrix is the ordered product
\begin{equation}
\label{eq:monodromy-matrix-preview}
    M_a(u)
    =
    R_{aL}(u-\xi_L)R_{a,L-1}(u-\xi_{L-1})\cdots R_{a1}(u-\xi_1),
\end{equation}
where $\xi_j$ are inhomogeneity parameters.  The homogeneous chain is obtained by
setting all $\xi_j=0$.  The transfer matrix is the auxiliary trace
\begin{equation}
\label{eq:transfer-matrix-from-monodromy-preview}
    \tau(u)=\Tr_a M_a(u).
\end{equation}

The Yang--Baxter equation implies the RTT relation
\begin{equation}
\label{eq:rtt-relation-preview}
    R_{ab}(u-v)M_a(u)M_b(v)
    =
    M_b(v)M_a(u)R_{ab}(u-v) .
\end{equation}
Taking the trace over the two auxiliary spaces $a$ and $b$ gives
\begin{equation}
\label{eq:commuting-transfer-matrices-preview}
    [\tau(u),\tau(v)]=0,
    \qquad
    \text{for all }u,v .
\end{equation}
This is the central algebraic output of Yang--Baxter integrability.

\begin{theorem}[Commuting transfer matrices]
If the $R$-matrix satisfies the Yang--Baxter equation and the monodromy matrix is
constructed as in \cref{eq:monodromy-matrix-preview}, then the transfer matrices
\cref{eq:transfer-matrix-from-monodromy-preview} commute for all values of the
spectral parameter:
\begin{equation}
    [\tau(u),\tau(v)]=0 .
\end{equation}
\end{theorem}

\begin{proof}
The Yang--Baxter equation implies the RTT relation
\cref{eq:rtt-relation-preview}.  Taking $\Tr_a\Tr_b$ of both sides and using
cyclicity of the trace in the auxiliary spaces gives
\begin{equation}
    \Tr_a\Tr_b\left[R_{ab}(u-v)M_a(u)M_b(v)\right]
    =
    \Tr_a\Tr_b\left[M_b(v)M_a(u)R_{ab}(u-v)\right].
\end{equation}
Cyclicity moves $R_{ab}(u-v)$ through the auxiliary trace, leaving
\begin{equation}
    \Tr_a M_a(u)\,\Tr_b M_b(v)
    =
    \Tr_b M_b(v)\,\Tr_a M_a(u).
\end{equation}
This is exactly $\tau(u)\tau(v)=\tau(v)\tau(u)$.
\end{proof}

The physical meaning of \cref{eq:commuting-transfer-matrices-preview} is that
$\tau(u)$ is a generating function for conserved quantities.  Expanding the
logarithm around a regular point gives
\begin{equation}
\label{eq:charges-from-log-transfer-preview}
    \log\tau(u)
    =
    \sum_{n=0}^{\infty} (u-u_0)^n Q_n,
    \qquad
    [Q_m,Q_n]=0 .
\end{equation}
One of these charges is the Hamiltonian; the rest are higher conserved charges.
This is the transfer-matrix form of quantum integrability.

\section{Hamiltonian limit}
\label{sec:hamiltonian-limit-preview}

A local quantum Hamiltonian is obtained from a regular transfer matrix.  The key
regularity property is
\begin{equation}
\label{eq:R-regularity-preview}
    R(0)=\rho P,
\end{equation}
where $P$ is the permutation operator on two neighboring spin spaces and $\rho$ is
a scalar normalization.  Then the logarithmic derivative of the transfer matrix
at $u=0$ produces a sum of local nearest-neighbor terms:
\begin{equation}
\label{eq:hamiltonian-log-transfer-preview}
    H
    =
    C\left.\frac{\dd}{\dd u}\log\tau(u)\right|_{u=0}
    +C_0
    =
    \sum_{j=1}^{L}h_{j,j+1}+C_0 .
\end{equation}
For the six-vertex $R$-matrix, this Hamiltonian is the spin-$1/2$ XXZ chain,
up to normalization, additive constants, and boundary conventions:
\begin{equation}
\label{eq:xxz-from-six-vertex-preview}
    H_{\rm XXZ}
    =
    J\sum_{j=1}^{L}
    \left(
    S_j^xS_{j+1}^x
    +S_j^yS_{j+1}^y
    +\Delta S_j^zS_{j+1}^z
    \right)
    +\text{constant} .
\end{equation}
This is the basic bridge between the six-vertex model and the XXZ quantum chain.
It is the interacting analogue of the quantum--classical correspondence discussed
in the transfer-matrix chapter.

\section{Why commuting transfer matrices imply many conserved quantities}
\label{sec:commuting-transfer-conserved-quantities}

Suppose $\tau(u)$ is analytic near $u_0$ and $[\tau(u),\tau(v)]=0$ for all $u,v$.
Then every coefficient in the expansion of $\log\tau(u)$ commutes with every other
coefficient.  More explicitly, define
\begin{equation}
    Q_n
    =
    \left.
    \frac{\dd^n}{\dd u^n}\log\tau(u)
    \right|_{u=u_0} .
\end{equation}
Since the whole family $\tau(u)$ commutes, one obtains
\begin{equation}
    [Q_m,Q_n]=0,
    \qquad
    [Q_n,H]=0 .
\end{equation}
The first nontrivial local charge is usually the Hamiltonian.  Higher charges
have increasing spatial range.  In a local integrable chain one may schematically
write
\begin{equation}
    Q_2=H,
    \qquad
    Q_3=\sum_j q_{j,j+1,j+2},
    \qquad
    Q_4=\sum_j q_{j,j+1,j+2,j+3},
    \quad \ldots
\end{equation}
The existence of these charges is the algebraic reason why scattering is
elastic and factorized.

\begin{remark}[Integrability is stronger than one conservation law]
Translation invariance gives momentum conservation, and time-translation
invariance gives energy conservation.  These two conservation laws are not enough
to solve an interacting many-body problem.  Bethe-ansatz integrability requires an
infinite tower of mutually commuting charges, generated by a commuting family of
transfer matrices.
\end{remark}

\section{Coordinate Bethe ansatz versus algebraic Bethe ansatz}
\label{sec:coordinate-vs-algebraic-ba-preview}

There are two complementary ways to use this structure.

\subsection{Coordinate Bethe ansatz}

The coordinate Bethe ansatz works directly with wave functions such as
\cref{eq:coordinate-ba-wavefunction-preview}.  Its main steps are:
\begin{enumerate}[label=(\roman*)]
    \item choose a simple reference state;
    \item construct $M$-particle wave functions in ordered coordinate sectors;
    \item impose contact conditions when particles become adjacent;
    \item extract two-body scattering phases;
    \item impose periodic boundary conditions to obtain Bethe equations.
\end{enumerate}
This method is concrete and physically transparent.  It is especially useful for
the XXX and XXZ chains and will be developed in the coordinate Bethe-ansatz
chapter.

\subsection{Algebraic Bethe ansatz}

The algebraic Bethe ansatz starts from the monodromy matrix.  For a spin-$1/2$
chain the monodromy matrix is a $2\times2$ matrix in the auxiliary space,
\begin{equation}
\label{eq:monodromy-ABCD-preview}
    M_a(u)
    =
    \begin{pmatrix}
    A(u)&B(u)\\
    C(u)&D(u)
    \end{pmatrix}_a,
\end{equation}
where $A,B,C,D$ are operators acting on the quantum spin chain.  The transfer
matrix is
\begin{equation}
    \tau(u)=A(u)+D(u).
\end{equation}
One chooses a reference state $\ket{0}$ such that
\begin{equation}
    C(u)\ket{0}=0,
    \qquad
    A(u)\ket{0}=a(u)\ket{0},
    \qquad
    D(u)\ket{0}=d(u)\ket{0} .
\end{equation}
The operator $B(u)$ creates Bethe excitations:
\begin{equation}
\label{eq:algebraic-bethe-state-preview}
    \ket{\lambda_1,\ldots,\lambda_M}
    =B(\lambda_1)\cdots B(\lambda_M)\ket{0} .
\end{equation}
The RTT relation determines how $A(u)$ and $D(u)$ commute through the $B$'s.  The
unwanted terms vanish exactly when the rapidities satisfy the Bethe equations.
Then \cref{eq:algebraic-bethe-state-preview} is an eigenstate of $\tau(u)$ and
therefore also of the Hamiltonian extracted from $\log\tau(u)$.

\begin{center}
\begin{tabularx}{\textwidth}{L{0.28\textwidth}Y Y}
\toprule
Approach & Main object & Best use \\
\midrule
Coordinate Bethe ansatz & Piecewise plane-wave many-body wave function & Direct derivation of scattering phases and Bethe equations \\
Algebraic Bethe ansatz & Monodromy matrix and RTT algebra & Systematic construction of eigenstates from $B(u)$ operators \\
Transfer-matrix method & Commuting family $\tau(u)$ & Conserved charges and Hamiltonian limits of vertex models \\
\bottomrule
\end{tabularx}
\end{center}

\section{Relation to free-fermion solvability}
\label{sec:bethe-vs-free-fermion-preview}

It is useful to compare the free-fermion and Bethe-ansatz mechanisms.

\begin{center}
\begin{tabularx}{\textwidth}{L{0.26\textwidth}Y Y}
\toprule
Feature & Free-fermion/BdG solvability & Bethe/Yang--Baxter integrability \\
\midrule
Elementary objects & Independent fermionic modes or quasiparticles & Interacting Bethe quasiparticles with rapidities \\
Many-body construction & Fill independent modes & Solve coupled Bethe equations \\
Scattering & Trivial exchange signs or BdG mixing & Nontrivial two-body phase or matrix scattering \\
Conserved quantities & Quadratic mode occupations & Infinite tower from $\log\tau(u)$ \\
Representative models & XY, TFIM, Kitaev, SSH & XXX, XXZ, XYZ, six-vertex, Hubbard \\
Core algebra & Canonical anticommutation and Bogoliubov rotation & Yang--Baxter equation and RTT relation \\
\bottomrule
\end{tabularx}
\end{center}

The XX chain sits at the overlap between these viewpoints.  It is the
$\Delta=0$ limit of the XXZ chain and is also free after the Jordan--Wigner
transformation.  From the Bethe-ansatz perspective, its scattering phase reduces
to the fermionic exchange sign.  From the free-fermion perspective, it is simply a
number-conserving tight-binding chain.  This overlap is a useful consistency
check between the two solution languages.

\section{Physical interpretation of rapidities}
\label{sec:rapidities-preview}

In relativistic field theory a rapidity directly parametrizes energy and momentum.
In lattice Bethe ansatz, the word is used more broadly: a rapidity is a spectral
parameter that uniformizes the dispersion and scattering.  The physical momentum
is a function of rapidity,
\begin{equation}
    p=p(\lambda),
\end{equation}
and the two-body scattering depends only on rapidity differences in many
homogeneous integrable models,
\begin{equation}
    S(\lambda_i,\lambda_j)=S(\lambda_i-\lambda_j).
\end{equation}
This difference property is one reason why the Yang--Baxter equation becomes a
tractable functional equation for $R(u)$.

For the XXX chain one often uses
\begin{equation}
\label{eq:xxx-momentum-rapidity-preview}
    \ee^{\ii p(\lambda)}
    =
    \frac{\lambda+\ii/2}{\lambda-\ii/2} .
\end{equation}
Then \cref{eq:xxx-bethe-equations-preview} has the direct interpretation
\begin{equation}
    \ee^{\ii p(\lambda_j)L}
    =
    \prod_{\ell\neq j}S(\lambda_j-\lambda_\ell),
    \qquad
    S(\lambda)=\frac{\lambda+\ii}{\lambda-\ii} .
\end{equation}

\section{What Yang--Baxter integrability does not mean}
\label{sec:what-integrability-does-not-mean}

The term ``integrable'' is used in several related but distinct senses.  In these
notes, Yang--Baxter integrability will mean the existence of an $R$-matrix
satisfying the Yang--Baxter equation, a commuting transfer-matrix family, and a
Hamiltonian obtained from a logarithmic derivative.  This is stronger than merely
knowing many exact eigenstates or finding one special solvable point.

There are several useful cautions.
\begin{enumerate}[label=(\roman*)]
    \item A model may be exactly solvable by free fermions without requiring a
    nontrivial Yang--Baxter structure.
    \item A model may have a transfer matrix but not a commuting family
    $\tau(u)$ for all spectral parameters.
    \item Self-duality may identify a candidate critical point, but it is not the
    same as Bethe-ansatz integrability.
    \item Topological two-band classification concerns eigenvector winding or
    Chern structure; Yang--Baxter integrability concerns factorized scattering
    and commuting transfer matrices.
\end{enumerate}
The value of the mother-model perspective is that these solution mechanisms can
be compared without conflating them.

\section{Roadmap to the later integrable chapters}
\label{sec:integrability-roadmap-preview}

The material in this preview will be used later in three directions.

\begin{enumerate}[label=(\roman*)]
    \item The six-vertex model will provide the canonical example of an
    $R$-matrix satisfying the Yang--Baxter equation.
    \item The Hamiltonian limit of the six-vertex transfer matrix will give the
    XXZ spin chain.
    \item The coordinate Bethe ansatz will solve the XXZ chain by deriving the
    many-body Bethe equations from factorized scattering.
    \item The algebraic Bethe ansatz will solve the same chain from the RTT
    relation and the operator algebra of $A(u),B(u),C(u),D(u)$.
    \item The Hubbard chain will require a nested Bethe ansatz because its
    quasiparticles carry more than one kind of internal degree of freedom.
\end{enumerate}

\section{Summary}
\label{sec:bethe-yang-baxter-summary}

The essential ideas are:
\begin{enumerate}[label=(\roman*)]
    \item Bethe ansatz replaces independent free modes by interacting rapidities
    whose many-body scattering factorizes into two-body scattering.
    \item The coordinate Bethe wave function is a sum of plane waves over
    momentum permutations, with amplitudes related by two-body scattering
    matrices or phases.
    \item Periodic boundary conditions give Bethe equations of the form
    $\ee^{\ii k_jL}=\prod_{\ell\neq j}S(k_j,k_\ell)$.
    \item For matrix-valued scattering, three-body consistency is the
    Yang--Baxter equation.
    \item The $R$-matrix is both a local vertex weight in a two-dimensional
    classical model and the algebraic scattering/intertwining object of a quantum
    chain.
    \item The Yang--Baxter equation implies the RTT relation, which implies
    $[\tau(u),\tau(v)]=0$.
    \item Expanding $\log\tau(u)$ produces mutually commuting conserved charges;
    the quantum Hamiltonian appears as one of these charges.
    \item The six-vertex model and the XXZ chain are the basic interacting
    examples connecting transfer matrices, Bethe ansatz, and Yang--Baxter
    integrability.
\end{enumerate}
